% Generated draft: requires mathematical and source review.
\documentclass[11pt]{article}
\usepackage[margin=25mm]{geometry}
\usepackage{fontspec}
\setmainfont{Latin Modern Roman}
\usepackage{amsmath,amssymb,mathtools,unicode-math}
\setmathfont{Latin Modern Math}
\usepackage{enumitem,needspace,xurl,hyperref}
\setlength{\parindent}{0pt}
\setlength{\parskip}{5pt}
\setlength{\emergencystretch}{3em}
\newcommand{\N}{\mathbb{N}}
\newcommand{\Z}{\mathbb{Z}}
\newcommand{\Q}{\mathbb{Q}}
\newcommand{\R}{\mathbb{R}}
\newcommand{\C}{\mathbb{C}}
\newcommand{\F}{\mathbb{F}}
\newcommand{\A}{\mathbb{A}}
\newcommand{\PP}{\mathbb{P}}
\newcommand{\E}{\mathbb{E}}
\newcommand{\eps}{\varepsilon}
\newcommand{\dd}{\,\mathrm{d}}
\begin{document}
\section*{UnsolvedMath ID: 701145}
\textbf{Source problem code:} catalogue entry
\subsection{Definitions and mathematical statement}
Let
\[
G=\operatorname{SL}_2(\Q)=\{A\in M_2(\Q):\det A=1\}.
\]
Equip \(G\) with the Zariski topology induced by the affine algebraic group \(\operatorname{SL}_2\) over \(\Q\). Thus a subgroup \(H\leq G\) is Zariski dense if every polynomial \(f\in\Q[x_{11},x_{12},x_{21},x_{22}]\) that vanishes on \(H\) also vanishes on all of \(G\). Say that \(H\) has no nontrivial finite quotients if, for every finite group \(F\) and every group homomorphism \(\varphi:H\to F\), one has \(\varphi(H)=\{1_F\}\).

Is it true that, for every subgroup \(H\leq G\),
\[
\bigl(H\text{ is Zariski dense in }G\ \text{and }H\text{ has no nontrivial finite quotients}\bigr)
\quad\Longrightarrow\quad H=G?
\]
\subsection{Short English statement}
Must every Zariski dense subgroup of \(\operatorname{SL}_2(\Q)\) with no nontrivial finite quotients equal \(\operatorname{SL}_2(\Q)\)? The supplied registry labels this problem open; that status is not independently verified.
\subsection{Sources}
\begin{itemize}
\item Local UnsolvedMath registry, supplied source record with permanent ID \(701145\), exact source code \texttt{catalogue entry}, and upstream revision \texttt{d371ee917ac9647fbb06cd6fc529deec6f3dda82}.
\item The supplied record attributes Problem 15.57 to J. McKay and J.-P. Serre, proposed in 2002, in \emph{The Kourovka Notebook (Unsolved Problems in Group Theory)}, 15th issue (2002). It also identifies the 21st issue edition (2026), \url{https://arxiv.org/abs/1401.0300v48}. The supplied collection description names E. I. Khukhro and V. D. Mazurov as editors. These bibliographic details are supplied attribution and have not been independently checked.
\end{itemize}
\end{document}
